\chapter{Low-Allocation Renderer}\label{record-052}

\noindent{\footnotesize\color{BookMuted}\path{docs/strategy/low-allocation-renderer.md}\par}\smallskip

\emph{Design note for reusable framebuffers and a low-garbage hot render path}

\section{Why this exists}\label{record-052-why-this-exists}

The DOGFOOD title screen made a real renderer truth visible:

\begin{itemize}
\tightlist
\item
  app-side title-screen work was slow enough to optimize
\item
  terminal paint cost still dominates at very large window sizes
\item
  even after the title-screen perf pass, the core runtime and surface stack still generate a lot of temporary garbage every frame
\end{itemize}

That means the next performance step is not another one-off screen optimization. It is renderer architecture.

The question is not:

\begin{quote}
``Can we make this one shader faster?''
\end{quote}

The real question is:

\begin{quote}
``Can the Bijou runtime render rich full-screen surfaces without allocating a storm of temporary objects every frame?''
\end{quote}

The answer is:

\begin{itemize}
\tightlist
\item
  \textbf{we can get much closer}
\item
  \textbf{but not with the current runtime/surface hot path}
\end{itemize}

\section{Evidence}\label{record-052-evidence}

\subsection{DOGFOOD title-screen benchmark}\label{record-052-dogfood-title-screen-benchmark}

At \texttt{2\allowbreak{}2\allowbreak{}0\allowbreak{}x\allowbreak{}5\allowbreak{}8}, the optimized DOGFOOD title screen now benchmarks at:

\begin{itemize}
\tightlist
\item
  \texttt{3\allowbreak{}.\allowbreak{}3\allowbreak{}9\allowbreak{} \allowbreak{}m\allowbreak{}s\allowbreak{}/\allowbreak{}f\allowbreak{}r\allowbreak{}a\allowbreak{}m\allowbreak{}e} for update + view + normalize
\item
  \texttt{3\allowbreak{}.\allowbreak{}9\allowbreak{}8\allowbreak{} \allowbreak{}m\allowbreak{}s\allowbreak{}/\allowbreak{}f\allowbreak{}r\allowbreak{}a\allowbreak{}m\allowbreak{}e} for update + view + normalize + diff/write into test IO
\end{itemize}

That is a major improvement over the earlier path:

\begin{itemize}
\tightlist
\item
  \texttt{1\allowbreak{}2\allowbreak{}.\allowbreak{}2\allowbreak{}7\allowbreak{} \allowbreak{}m\allowbreak{}s\allowbreak{}/\allowbreak{}f\allowbreak{}r\allowbreak{}a\allowbreak{}m\allowbreak{}e} before for update + view + normalize
\item
  \texttt{1\allowbreak{}1\allowbreak{}.\allowbreak{}8\allowbreak{}1\allowbreak{} \allowbreak{}m\allowbreak{}s\allowbreak{}/\allowbreak{}f\allowbreak{}r\allowbreak{}a\allowbreak{}m\allowbreak{}e} before for update + view + normalize + diff/write
\end{itemize}

So there was real waste in the title screen itself, and it was worth removing.

\subsection{Allocation profile}\label{record-052-allocation-profile}

The remaining problem is not retained memory. It is temporary churn.

Measured over \texttt{6\allowbreak{}0\allowbreak{}0} DOGFOOD landing frames at \texttt{2\allowbreak{}2\allowbreak{}0\allowbreak{}x\allowbreak{}5\allowbreak{}8}:

\begin{itemize}
\tightlist
\item
  transient heap delta: about \texttt{6\allowbreak{}.\allowbreak{}6\allowbreak{} \allowbreak{}M\allowbreak{}B}
\item
  retained heap delta after forced GC: about \texttt{2\allowbreak{}1\allowbreak{}3\allowbreak{} \allowbreak{}K\allowbreak{}B}
\end{itemize}

So this is \textbf{not a leak}. It is \textbf{temporary garbage production}.

\subsection{GC behavior}\label{record-052-gc-behavior}

The trace confirms that the runtime is feeding the collector frequently:

\begin{itemize}
\tightlist
\item
  repeated minor scavenges
\item
  at least one old-space mark-compact pause around \texttt{2\allowbreak{}4\allowbreak{} \allowbreak{}m\allowbreak{}s}
\end{itemize}

That is enough to matter for animation smoothness and pacing.

\section{The current bottlenecks}\label{record-052-the-current-bottlenecks}

\subsection{\texorpdfstring{1. \texttt{S\allowbreak{}u\allowbreak{}r\allowbreak{}f\allowbreak{}a\allowbreak{}c\allowbreak{}e} is object-heavy by design}{1.  is object-heavy by design}}\label{record-052-surface-is-object-heavy-by-design}

The current \texttt{S\allowbreak{}u\allowbreak{}r\allowbreak{}f\allowbreak{}a\allowbreak{}c\allowbreak{}e} API is ergonomic, but not cheap:

\begin{itemize}
\tightlist
\item
  \texttt{g\allowbreak{}e\allowbreak{}t\allowbreak{}(\allowbreak{})} returns masked cell copies
\item
  \texttt{s\allowbreak{}e\allowbreak{}t\allowbreak{}(\allowbreak{})} writes through \texttt{a\allowbreak{}p\allowbreak{}p\allowbreak{}l\allowbreak{}y\allowbreak{}M\allowbreak{}a\allowbreak{}s\allowbreak{}k\allowbreak{}(\allowbreak{}.\allowbreak{}.\allowbreak{}.\allowbreak{})}, which returns a fresh \texttt{C\allowbreak{}e\allowbreak{}l\allowbreak{}l}
\item
  \texttt{c\allowbreak{}l\allowbreak{}e\allowbreak{}a\allowbreak{}r\allowbreak{}(\allowbreak{})} replaces every cell object
\item
  \texttt{c\allowbreak{}l\allowbreak{}o\allowbreak{}n\allowbreak{}e\allowbreak{}(\allowbreak{})} deep-copies the whole surface
\item
  \texttt{g\allowbreak{}e\allowbreak{}t\allowbreak{}R\allowbreak{}o\allowbreak{}w\allowbreak{}(\allowbreak{})} allocates both arrays and masked cell copies
\end{itemize}

This is fine for a friendly public API. It is not fine for a hot renderer.

\subsection{2. The runtime allocates a fresh target surface every frame}\label{record-052-the-runtime-allocates-a-fresh-target-surface-every-frame}

Today the runtime:

\begin{itemize}
\tightlist
\item
  allocates a new \texttt{t\allowbreak{}a\allowbreak{}r\allowbreak{}g\allowbreak{}e\allowbreak{}t\allowbreak{}S\allowbreak{}u\allowbreak{}r\allowbreak{}f\allowbreak{}a\allowbreak{}c\allowbreak{}e} for every render
\item
  diffs it against \texttt{c\allowbreak{}u\allowbreak{}r\allowbreak{}r\allowbreak{}e\allowbreak{}n\allowbreak{}t\allowbreak{}S\allowbreak{}u\allowbreak{}r\allowbreak{}f\allowbreak{}a\allowbreak{}c\allowbreak{}e}
\item
  then clones the target into \texttt{c\allowbreak{}u\allowbreak{}r\allowbreak{}r\allowbreak{}e\allowbreak{}n\allowbreak{}t\allowbreak{}S\allowbreak{}u\allowbreak{}r\allowbreak{}f\allowbreak{}a\allowbreak{}c\allowbreak{}e}
\end{itemize}

That means the frame loop is paying for:

\begin{itemize}
\tightlist
\item
  one full-surface allocation
\item
  one full-surface clone
\end{itemize}

on every frame, even before terminal paint cost is considered.

\subsection{3. The diff path allocates while reading cells}\label{record-052-the-diff-path-allocates-while-reading-cells}

\texttt{r\allowbreak{}e\allowbreak{}n\allowbreak{}d\allowbreak{}e\allowbreak{}r\allowbreak{}D\allowbreak{}i\allowbreak{}f\allowbreak{}f\allowbreak{}(\allowbreak{}.\allowbreak{}.\allowbreak{}.\allowbreak{})} currently walks the grid through \texttt{s\allowbreak{}u\allowbreak{}r\allowbreak{}f\allowbreak{}a\allowbreak{}c\allowbreak{}e\allowbreak{}.\allowbreak{}g\allowbreak{}e\allowbreak{}t\allowbreak{}(\allowbreak{}.\allowbreak{}.\allowbreak{}.\allowbreak{})}.

That means the diff hot loop allocates just to inspect cells. It also:

\begin{itemize}
\tightlist
\item
  builds temporary token objects
\item
  builds \texttt{b\allowbreak{}a\allowbreak{}t\allowbreak{}c\allowbreak{}h\allowbreak{}T\allowbreak{}e\allowbreak{}x\allowbreak{}t} strings
\item
  builds a growing \texttt{o\allowbreak{}u\allowbreak{}t\allowbreak{}p\allowbreak{}u\allowbreak{}t} string
\end{itemize}

Some string allocation is unavoidable in an ANSI renderer. The per-cell object churn is not.

\subsection{4. The style layer is string-oriented}\label{record-052-the-style-layer-is-string-oriented}

The Node style adapter ultimately produces styled strings. That is expected, but it means the final bridge to the terminal will never be literally zero-allocation.

So the goal should be:

\begin{itemize}
\tightlist
\item
  \textbf{near-zero-allocation in the hot render loop}
\item
  not a fantasy of zero allocation everywhere
\end{itemize}

\section{The decision}\label{record-052-the-decision}

Bijou should pursue a \textbf{low-allocation renderer}, not promise a literally zero-allocation runtime.

The design decision is:

\begin{enumerate}
\def\labelenumi{\arabic{enumi}.}
\tightlist
\item
  keep the current \texttt{S\allowbreak{}u\allowbreak{}r\allowbreak{}f\allowbreak{}a\allowbreak{}c\allowbreak{}e} API as the ergonomic public composition model
\item
  stop treating that API as the final hot-path framebuffer contract
\item
  introduce reusable framebuffers and direct-cell diffing in the runtime
\item
  move toward a dedicated internal mutable render buffer if the first stage proves worthwhile
\end{enumerate}

In short:

\begin{quote}
\texttt{S\allowbreak{}u\allowbreak{}r\allowbreak{}f\allowbreak{}a\allowbreak{}c\allowbreak{}e} remains the authoring/composition type.\\
The runtime grows a lower-allocation framebuffer path underneath it.
\end{quote}

\section{Non-goals}\label{record-052-non-goals}

This effort does \textbf{not} mean:

\begin{itemize}
\tightlist
\item
  rewriting the whole design system around typed arrays immediately
\item
  removing all object allocation from every part of Bijou
\item
  breaking the public \texttt{S\allowbreak{}u\allowbreak{}r\allowbreak{}f\allowbreak{}a\allowbreak{}c\allowbreak{}e} API in the first slice
\item
  optimizing away terminal-emulator paint costs that live outside Bijou
\end{itemize}

Those would blur the work and make it much riskier than it needs to be.

\section{Principles}\label{record-052-principles}

\begin{enumerate}
\def\labelenumi{\arabic{enumi}.}
\item
  \textbf{Fix the hot path first} We should remove runtime/framebuffer churn before inventing a whole new render substrate.
\item
  \textbf{Preserve public ergonomics} App authors should not need to think in framebuffers just to build normal components.
\item
  \textbf{Make each step benchmarkable} Every stage should have before/after benchmarks and GC observations.
\item
  \textbf{Tests are the spec} Runtime, diff, and resize behavior must stay covered while the internals change.
\item
  \textbf{Avoid premature big-bang rewrites} The first slice should be a contained runtime improvement, not a wholesale new rendering engine.
\end{enumerate}

\section{Proposed architecture}\label{record-052-proposed-architecture}

\subsection{Stage 1 --- Runtime framebuffer reuse}\label{record-052-stage-1-runtime-framebuffer-reuse}

Replace the current per-frame target allocation and post-diff clone with a double-buffer model:

\begin{itemize}
\tightlist
\item
  allocate \texttt{f\allowbreak{}r\allowbreak{}o\allowbreak{}n\allowbreak{}t\allowbreak{}B\allowbreak{}u\allowbreak{}f\allowbreak{}f\allowbreak{}e\allowbreak{}r} once per viewport size
\item
  allocate \texttt{b\allowbreak{}a\allowbreak{}c\allowbreak{}k\allowbreak{}B\allowbreak{}u\allowbreak{}f\allowbreak{}f\allowbreak{}e\allowbreak{}r} once per viewport size
\item
  render into the back buffer
\item
  diff front vs back
\item
  swap references
\end{itemize}

This removes the largest guaranteed per-frame allocations without changing public component APIs.

\subsection{Stage 2 --- Direct-cell diff fast path}\label{record-052-stage-2-direct-cell-diff-fast-path}

Teach \texttt{r\allowbreak{}e\allowbreak{}n\allowbreak{}d\allowbreak{}e\allowbreak{}r\allowbreak{}D\allowbreak{}i\allowbreak{}f\allowbreak{}f\allowbreak{}(\allowbreak{}.\allowbreak{}.\allowbreak{}.\allowbreak{})} to read cells directly from \texttt{s\allowbreak{}u\allowbreak{}r\allowbreak{}f\allowbreak{}a\allowbreak{}c\allowbreak{}e\allowbreak{}.\allowbreak{}c\allowbreak{}e\allowbreak{}l\allowbreak{}l\allowbreak{}s} in the hot loop instead of calling \texttt{s\allowbreak{}u\allowbreak{}r\allowbreak{}f\allowbreak{}a\allowbreak{}c\allowbreak{}e\allowbreak{}.\allowbreak{}g\allowbreak{}e\allowbreak{}t\allowbreak{}(\allowbreak{}.\allowbreak{}.\allowbreak{}.\allowbreak{})}.

That should eliminate:

\begin{itemize}
\tightlist
\item
  masked read allocations
\item
  extra per-cell object copying during diff
\end{itemize}

It also creates a clearer path to style-run caching later.

\subsection{Stage 3 --- Mutable/internal render buffer}\label{record-052-stage-3-mutableinternal-render-buffer}

If stages 1 and 2 are not enough, add an internal runtime type such as:

\begin{itemize}
\tightlist
\item
  \texttt{F\allowbreak{}r\allowbreak{}a\allowbreak{}m\allowbreak{}e\allowbreak{}B\allowbreak{}u\allowbreak{}f\allowbreak{}f\allowbreak{}e\allowbreak{}r}
\item
  \texttt{M\allowbreak{}u\allowbreak{}t\allowbreak{}a\allowbreak{}b\allowbreak{}l\allowbreak{}e\allowbreak{}S\allowbreak{}u\allowbreak{}r\allowbreak{}f\allowbreak{}a\allowbreak{}c\allowbreak{}e}
\item
  or \texttt{R\allowbreak{}e\allowbreak{}n\allowbreak{}d\allowbreak{}e\allowbreak{}r\allowbreak{}S\allowbreak{}u\allowbreak{}r\allowbreak{}f\allowbreak{}a\allowbreak{}c\allowbreak{}e}
\end{itemize}

This type should:

\begin{itemize}
\tightlist
\item
  mutate cells in place
\item
  support clear/fill/blit without per-cell object replacement
\item
  expose direct array access for diffing
\end{itemize}

This is the likely bridge between the public \texttt{S\allowbreak{}u\allowbreak{}r\allowbreak{}f\allowbreak{}a\allowbreak{}c\allowbreak{}e} model and a truly lean renderer.

\subsection{Stage 4 --- Packed style and glyph data}\label{record-052-stage-4-packed-style-and-glyph-data}

Longer-term, the renderer can move toward:

\begin{itemize}
\tightlist
\item
  compact style ids
\item
  bitflags for modifiers
\item
  glyph/string lookup tables
\item
  packed arrays instead of per-cell object shapes
\end{itemize}

This is where the renderer starts behaving more like a real framebuffer instead of a grid of JS objects.

\subsection{Stage 5 --- ANSI emission optimization}\label{record-052-stage-5-ansi-emission-optimization}

Once the buffer model is leaner, improve the terminal bridge itself:

\begin{itemize}
\tightlist
\item
  cache style runs
\item
  reduce repeated SGR construction
\item
  avoid unnecessary string concatenation churn
\item
  write larger coherent chunks
\end{itemize}

That still will not be literal zero-allocation, but it will be much closer to the practical lower bound for ANSI output.

\section{Scope for the next implementation slice}\label{record-052-scope-for-the-next-implementation-slice}

\subsection{In scope}\label{record-052-in-scope}

\begin{itemize}
\tightlist
\item
  reusable front/back surfaces in the runtime
\item
  no per-frame \texttt{t\allowbreak{}a\allowbreak{}r\allowbreak{}g\allowbreak{}e\allowbreak{}t\allowbreak{}S\allowbreak{}u\allowbreak{}r\allowbreak{}f\allowbreak{}a\allowbreak{}c\allowbreak{}e\allowbreak{} \allowbreak{}=\allowbreak{} \allowbreak{}c\allowbreak{}r\allowbreak{}e\allowbreak{}a\allowbreak{}t\allowbreak{}e\allowbreak{}S\allowbreak{}u\allowbreak{}r\allowbreak{}f\allowbreak{}a\allowbreak{}c\allowbreak{}e\allowbreak{}(\allowbreak{}.\allowbreak{}.\allowbreak{}.\allowbreak{})}
\item
  no per-frame \texttt{c\allowbreak{}u\allowbreak{}r\allowbreak{}r\allowbreak{}e\allowbreak{}n\allowbreak{}t\allowbreak{}S\allowbreak{}u\allowbreak{}r\allowbreak{}f\allowbreak{}a\allowbreak{}c\allowbreak{}e\allowbreak{} \allowbreak{}=\allowbreak{} \allowbreak{}t\allowbreak{}a\allowbreak{}r\allowbreak{}g\allowbreak{}e\allowbreak{}t\allowbreak{}S\allowbreak{}u\allowbreak{}r\allowbreak{}f\allowbreak{}a\allowbreak{}c\allowbreak{}e\allowbreak{}.\allowbreak{}c\allowbreak{}l\allowbreak{}o\allowbreak{}n\allowbreak{}e\allowbreak{}(\allowbreak{})}
\item
  direct-buffer swap on successful output
\item
  benchmarks and runtime tests proving equivalence
\end{itemize}

\subsection{Out of scope}\label{record-052-out-of-scope}

\begin{itemize}
\tightlist
\item
  a new public \texttt{S\allowbreak{}u\allowbreak{}r\allowbreak{}f\allowbreak{}a\allowbreak{}c\allowbreak{}e} API
\item
  packed typed-array cells
\item
  changing component authoring patterns
\item
  redesigning the style port
\item
  terminal-specific output throttling or adaptive quality logic
\end{itemize}

\section{Acceptance criteria}\label{record-052-acceptance-criteria}

The first slice should be considered successful if:

\begin{enumerate}
\def\labelenumi{\arabic{enumi}.}
\tightlist
\item
  runtime behavior is unchanged from the app author's point of view
\item
  framebuffers are reused across steady-state renders
\item
  the per-frame allocation profile improves materially
\item
  render benchmarks and smoke suites stay green
\item
  resize behavior still replaces buffers correctly when dimensions change
\end{enumerate}

\section{Open questions}\label{record-052-open-questions}

\begin{enumerate}
\def\labelenumi{\arabic{enumi}.}
\tightlist
\item
  Should the reusable back buffer be fully cleared each frame, or should some middleware rely on an already-empty surface contract that needs stronger enforcement?
\item
  Should the diff fast path be added to the existing \texttt{S\allowbreak{}u\allowbreak{}r\allowbreak{}f\allowbreak{}a\allowbreak{}c\allowbreak{}e} implementation, or should it be introduced only after a dedicated internal framebuffer exists?
\item
  How far can we push runtime reuse before the \texttt{S\allowbreak{}u\allowbreak{}r\allowbreak{}f\allowbreak{}a\allowbreak{}c\allowbreak{}e} object model itself becomes the dominant remaining cost?
\end{enumerate}

\section{Recommendation}\label{record-052-recommendation}

Proceed immediately with \textbf{Stage 1: runtime framebuffer reuse}.

It is the highest-leverage improvement because:

\begin{itemize}
\tightlist
\item
  it removes unavoidable per-frame allocations already visible in the runtime
\item
  it does not require a public API break
\item
  it sets up every later optimization
\item
  it is easy to validate with benchmarks and tests
\end{itemize}

After that, take \textbf{Stage 2: direct-cell diff fast path} before deciding whether Stage 3 needs a new internal buffer type.

